\section{Single-profile state reductions}\label{rec:section}

We now use the general angular estimates to bound positive counts of
pairs with nearly equal distances.  Fix one regularized probability
$\rho$, with profile $m$, and write $A_n,M(n),I(a,b),k(a,b)$ for its
profile data.  Both endpoint laws in an angular state have this profile;
their cells and bounded spatial multipliers may differ.

There are two kinds of recursive state.  A positive state records a
count of pairs with nearly equal distances and uses the tests scheduled by its current
interval.  An angular state uses the independent energy of
Definition~\ref{osc:states}, together with a retained base interval that
specifies which tests must still be present.  We define these data before
stating the reductions that preserve them.

\subsection{The two recursive states}

\begin{definition}[Positive distance state]\label{osc:positive-count}
Let $0\le a\le t\le N$ be integer depths, let
$i\in\{0,\ldots,J_0\}$, and let $X,Y$ be a separated pair of occupied
depth-$a$ cells of $\rho$, in the geometry of
Equation~\eqref{osc:separation}.  Define the set of pairs satisfying the tests
\[
 \mathcal P_i^{X,Y}(a,t)=
 \bigl\{(x,y)\in X\times Y:
 e(x,y)\in G_{i,\mathfrak t}(x)\cap G_{i,\mathfrak t}(y)
 \text{ for every }\mathfrak t\in\mathcal M(a,t)\bigr\},
\]
and let
\[
 \theta_i^{X,Y}(a,t)
 =\mathbf1_{\mathcal P_i^{X,Y}(a,t)}
       (\rho|_X\otimes\rho|_Y).
\]
This is an unnormalized pair restriction.  Put $C_0=100(d+1)$ and
\begin{equation}\label{osc:count}
 \mathsf D_i^{X,Y}(a,t)
 =2^{t-a}M(a)^{-4}
   \int\mathbf1_{\{|D(x,y)-D(x',y')|\le C_0 2^{-t}\}}
   \,d\theta_i^{X,Y}(a,t)(x,y)
   \,d\theta_i^{X,Y}(a,t)(x',y').
\end{equation}
A positive state consists of these depths, index and cells, with value
$\mathsf D_i^{X,Y}(a,t)$.  Its current interval is $[a,t]$; it has no
separate retained base.  We suppress the cells in $\mathsf D_i(a,t)$ when
they are fixed or are varied in a stated maximum.
\end{definition}

\begin{definition}[Recursive angular state]\label{rec:angular-state}
Let $\mathcal B=[a_0,t_0]\subseteq[0,N]$ be an integer interval, called
the retained base, and let $[a,v]$ be the current interval, with
\begin{equation}\label{osc:base-contract}
 0\le a_0\le a\le v\le t_0\le N,
 \qquad t_0\le v+\sigma N.
\end{equation}
Write
\[
 \mathcal M_{>a}(\mathcal B)=
 \{\mathfrak t\in\mathcal M(a_0,t_0):
                 \text{the anchor of }\mathfrak t\text{ is deeper than }a\}.
\]
A recursive angular state consists of these intervals, an index
$i\in\{0,\ldots,J_0-1\}$, two occupied depth-$a$ cells $X,Y$ of
$\rho$, and two admissible endpoint symbols.  No separation of the cells
is required.  Each symbol's descriptor list consists of exactly one occurrence of each descriptor in
$\mathcal M_{>a}(\mathcal B)$, at the common index $i$.
Lemma~\ref{test:descriptor-properties} gives a common fixed polynomial
size bound for these lists throughout the recursive family.  The state also
specifies a fixed derivative budget and common numerical bounds for its
smooth angular factors, as in Definition~\ref{osc:symbols}; these bounds
are independent of $N$.

The value of this state is the already defined general energy for those
data:
\[
 \mathsf H_i^{\mathcal B}(a,v):=\mathsf H^{X,Y}(a,v).
\]
The compact notation displays the current interval, base and index.  The
cells, symbols and fixed bounds remain part of the state data even when
they are not displayed.  We use a maximum only for finite choices of cells,
and a supremum whenever admissible symbols vary.  Neither operation varies
the displayed base or index, or the fixed bounds, unless this is stated.
\end{definition}

A descriptor's origin is not a further factor in a state.  For each
membership $\mathfrak t\in\mathcal M_{>a}(\mathcal B)$,
Lemma~\ref{test:descriptor-properties} supplies an originating interval
inside this particular base.  If the descriptor has anchor $p$ and
length $L$, that origin contains $[p-L,p+L]$ and supplies the parallel
descriptors $[g,p]$ for $p-L\le g<p$.  An origin can be chosen when it is
used; it need not be the current interval, and a deduplicated descriptor
need not have one origin valid for every list containing it.

The general graph estimate returns energies with the restricted old
symbols, including reached factors.  Suppose a one-profile graph child at
depth $a'$ carries the list $\mathcal M_{>a}(\mathcal B)$ from a prior
start $a\le a'$.  The removal operation after
Equation~\eqref{osc:angular-cs} deletes its factors anchored at depths at
most $a'$ at bounded cost, leaving exactly
$\mathcal M_{>a'}(\mathcal B)$.  Once the interval inequalities in
Equation~\eqref{osc:base-contract} also hold, the resulting general energy
is the value of a recursive angular state.

The identity responsible for the same normalization in every dimension is
\begin{equation}\label{rec:normalization}
 2^{d(b-a)}\left(\frac{M(b)}{M(a)}\right)^2
 =2^{d(b-a)}R^{-2I(a,b)}
 =R^{-2(A_b-A_a)}.
\end{equation}
Here the last equality uses precisely $S=d/2$.

\subsection{Positive collision estimates}

\begin{lemma}[Binning positive measures]\label{rec:binning}
Let \(\eta,\zeta\) be finite positive measures on the line and let
\(\varepsilon>0\).  Write
\[
 B_\varepsilon(\eta)
 =\sum_{j\in\mathbb Z}\eta([j\varepsilon,(j+1)\varepsilon))^2.
\]
For any fixed \(C\ge1\),
\begin{align}
 B_\varepsilon(\eta)
 &\le(\eta\otimes\eta)\{|s-s'|\le\varepsilon\},
 \label{rec:bins-lower}\\
 \sup_{u\in\mathbb R}
 (\eta\otimes\eta)\{|s-s'-u|\le C\varepsilon\}
 &\le C' B_\varepsilon(\eta),
 \label{rec:bins-shift}\\
 \iint(1+|s-s'|/\varepsilon)^{-2}\,d\eta(s)d\zeta(s')
 &\le C' B_\varepsilon(\eta)^{1/2}
                    B_\varepsilon(\zeta)^{1/2}.
 \label{rec:bins-kernel}
\end{align}
Moreover,
\begin{equation}\label{rec:collision-fourier}
 (\eta\otimes\eta)\{|s-s'|\le C\varepsilon\}
 \le C'\varepsilon
       \int_{|r|\le\varepsilon^{-1}}|\widehat\eta(r)|^2\,dr.
\end{equation}
\end{lemma}

\begin{proof}
Put \(b_j=\eta([j\varepsilon,(j+1)\varepsilon))\) and define \(c_j\)
using \(\zeta\).  Cauchy--Schwarz gives, for every integer \(h\),
\[
 \sum_jb_jb_{j+h}\le\sum_jb_j^2,
 \qquad
 \sum_jb_jc_{j+h}\le
       \left(\sum_jb_j^2\sum_jc_j^2\right)^{1/2}.
\]
An interval of fixed length \(2C\varepsilon\), with any center, permits
only a bounded number of bin offsets.  This proves
\eqref{rec:bins-lower}--\eqref{rec:bins-shift}.  The kernel in
\eqref{rec:bins-kernel} is bounded on the bin pair with offset \(h\) by
\(C(1+|h|)^{-2}\), whose sum is finite.

For the last assertion, choose a fixed sufficiently small \(b>0\) and
put \(K(t)=(\sin(bt)/(bt))^2\), with its continuous value at zero.
Then \(K\ge c>0\) on \([-C,C]\), and its nonnegative Fourier transform
is bounded and supported in \([-1,1]\).  Fourier inversion and Fubini,
applied to \(K((s-s')/\varepsilon)\), prove
\eqref{rec:collision-fourier}.  These arguments require no density for
either measure.
\end{proof}

\subsection{The reduction mechanism}

A positive state uses the fixed policy choice $p$ of its current interval
$[a,t]$.  At a suffix choice, the projection test for the full conditional
law permits a direct passage to the coarser distance count at depth $p$.
At a prefix choice, a positive Fourier estimate splits the count into
frequency shells.  Low shells give shorter positive states; high shells
use the weighted shell conversion and the graph estimate to enter an
angular state.

For an angular state, let $L$ be the largest length in its remaining
list, with $L=0$ for an empty list, and put
\begin{equation}\label{rec:angular-split-depth}
 L_0=\max\{L,\lceil10\sigma N\rceil\},\qquad h=v-2L_0.
\end{equation}
This compares the longest remaining test with half the current frequency
range.  If $h\le a$, an origin of a longest descriptor supplies a parallel
test for a graph step.  If $h>a$, splitting at depth $h$ leaves enough
frequency range for stationary phase on the far terms; the near terms
have only boundedly many neighboring cells.  Thus an angular episode
continues through direct or near steps and returns to a positive state
through a far step.

The current interval is $[a,t]$ for $\mathsf D_i(a,t)$ and $[a,v]$ for
$\mathsf H_i^{\mathcal B}(a,v)$.  Terminality refers to this interval:
its width is at most $100q_0$.  It does not refer to the retained base.

\begin{proposition}[State reductions]\label{rec:reductions}
Suppose the current state is nonterminal, its smoothing index satisfies
$i<J_0$, and $N$ is sufficiently large with all parameters fixed.  For
every prescribed $T>0$, it is bounded, apart from an error $O(R^{-T})$
from oscillatory operations, by a sum of child bounds whose number is
polynomial in $N$, and possibly a low-frequency stopping term.  For every
non-low child,
if the old start is $a$ and the new start is $a'$, its multiplier is
\begin{equation}\label{rec:transition-factor}
 R^{-2(A_{a'}-A_a)+\mathrm{cost}+C_d\sigma+o(1)}.
\end{equation}
A low child has cost zero and the entire multiplier $R^{o(1)}$, with no
fixed $C_d\sigma$ loss.  A low-frequency stopping term is
$R^{C_d\sigma+o(1)}$.  The legal children are
\begin{center}
\renewcommand{\arraystretch}{1.2}
\begin{tabular}{lll}
Reduction & Parent and child & Cost \\ \hline
Suffix & $\mathsf D_i(a,t)\to\mathsf D_{i+1}(a,p)$
 & $A_p-\min_{[p,t]}A$ \\
Low & $\mathsf D_i(a,t)\to\mathsf D_{i+1}(a,v)$
 & $0$; full factor $R^{o(1)}$ \\
High & $\mathsf D_i(a,t)\to\mathsf H_i^{[a,t]}(p,v)$
 & $k(a,p)$ \\
Direct & $\mathsf H_i^{\mathcal B}(a,v)\to\mathsf H_i^{\mathcal B}(p,v)$
 & $k(a,p)$ \\
Near & $\mathsf H_i^{\mathcal B}(a,v)\to\mathsf H_i^{\mathcal B}(h,v)$
 & $0$ \\
Far & $\mathsf H_i^{\mathcal B}(a,v)\to\mathsf D_{i+1}(b,v)$
 & $0$
\end{tabular}
\end{center}

For a positive parent, $p$ is its current interval's fixed policy choice.
The suffix row is used when $2p>a+t$.  When $2p\le a+t$, the shell index
is an integer $a<v\le t$: the low row has $t-v\ge\sigma N$, and the high
row has $0\le t-v<\sigma N$.  A high child has base $[a,t]$ and carries
the fixed enlarged derivative budget and common angular bounds produced
by the inverse conversion; these depend on $T$ and the fixed parameters,
not on $N$.

For an angular parent, use $L_0,h$ from
Equation~\eqref{rec:angular-split-depth}.  If $h\le a$, the direct row
chooses a remaining descriptor of maximal length $L$, its anchor $p$,
and an origin $[b_0,u_0]\subseteq\mathcal B$ for that membership.  This
choice has $L=L_0$, $a<p<v$,
\[
 b_0\le p-L,\qquad p+L\le u_0,\qquad v-a\le2L,
\]
and it supplies the parallel descriptor $[\max(a,p-L),p]$ used by the
graph estimate.  If $h>a$, the near row uses the displayed depth $h$.
The far rows range over integer depths $b$ arising from the comparable
far displacements, with $a<b\le h+O_d(1)<v$, and use separated depth-$b$
cell pairs.  Direct and near children keep the same base, index,
derivative budget and common angular bounds.  A far child has the fresh
positive list $\mathcal M(b,v)$.

Finite maxima over child cells use the stated geometry.  Whenever admissible
symbols vary, take a supremum over the symbol family just specified.  An
empty collection of child terms contributes zero.  Constants outside powers
of $R$ may depend on
the fixed parameters, derivative budgets and angular bounds.  The
constants $C_d$ in exponent losses depend only on $d$.  A terminal state
is bounded by
\begin{equation}\label{rec:terminal}
 R^{C_d(\kappa+\sigma)+o(1)}.
\end{equation}
\end{proposition}

\begin{proof}
We prove every reduction, including the support conditions needed for
its iteration.

\paragraph{Suffix.}
Let \(\theta=\theta_i^{X,Y}(a,t)\).  Retain a pair \((P,Q)\) of
depth-\(p\) subcells of the separated parent cells whenever
\(\mathcal P_i^{X,Y}(a,t)\cap(P\times Q)\ne\varnothing\), and choose
one pair in that intersection as a witness.  Put \(e_0=e(c_P,c_Q)\).  Separation of the parent cells gives
the uniform Hessian bound \(\|\nabla^2D\|\lesssim2^a\), and hence
\begin{equation}\label{rec:suffix-taylor}
 D(x,y)=e_0\cdot(x-y)+O_d(2^{a-2p})
       =e_0\cdot(x-y)+O_d(2^{-t})
       \quad(x\in P,\ y\in Q).
\end{equation}
The last equality follows from the suffix inequality \(2p>a+t\).
The witness angle differs from \(e_0\) by \(O_d(2^{a-p})\).
For \(y,y'\in Q\), changing between these two angles changes the
projection of \(y-y'\) by \(O_d(2^{a-2p})\).

The witness satisfies the projection test \([p,t]\).  This test is
formulated for the full conditional probability \(\rho|_Q/\rho(Q)\),
not for a witness-dependent restriction.  Its width absorbs every
fixed multiple of \(2^{-t}\).  Thus the conditional projected measure
in direction \(e_0\) has self-collision mass at tolerance
\(C2^{-t}\) bounded by
\begin{equation}\label{rec:suffix-B}
 B=2^{-(t-p)}R^{A_p-\min_{p\le n\le t}A_n+C_d\sigma}.
\end{equation}
By \eqref{rec:bins-shift}, the same bound, up to a fixed constant,
holds uniformly for any translated collision window.  Fixing
\(x,x'\in P\) and using \eqref{rec:suffix-taylor} therefore gives
\[
 B_{2^{-t}}\bigl(D_\#(\rho|_P\otimes\rho|_Q)\bigr)
 \lesssim B\rho(P)^2\rho(Q)^2.
\]

Group the retained cell pairs by the bin of width \(2^{-p}\)
containing their center distance, and write
\(W_G=\sum_{(P,Q)\in G}\rho(P)\rho(Q)\).
The fine distance-bin vector of a group is a sum of positive vectors.
The triangle inequality in \(\ell^2\) bounds its squared norm by
\(CBW_G^2\).  Each group's distances occupy a fixed enlargement of
its coarse bin, and these enlargements have bounded overlap.
Consequently the parent's fine collision mass is at most
\(CB\sum_GW_G^2\), by Lemma~\ref{rec:binning}.

Every full pair in a retained \(P\times Q\) satisfies the child list
\(\mathcal M(a,p)\) at level \(i+1\).  Indeed its descriptors belong
to the parent list, have anchors no deeper than \(p\), and have
nominal lengths at most \((p-a)/2\).  The spatial anchor cells of the
full pair and witness agree.  The change of direction is
\(O_d(2^{a-p})\), which lies inside the transfer neighborhood of
Lemma~\ref{test:smooth-masks} for every such descriptor.  Two full
pairs in one coarse group have distance difference at most
\((1+4\sqrt d)2^{-p}<C_0 2^{-p}\).  The Definition of the positive
child count now gives
\[
 \sum_GW_G^2
 \le 2^{-(p-a)}M(a)^4\mathsf D_{i+1}(a,p).
\]
Multiplying by the parent normalization, and substituting
\eqref{rec:suffix-B}, yields
\[
 2^{t-a}M(a)^{-4}\,B\,
 2^{-(p-a)}M(a)^4
 =R^{A_p-\min_{[p,t]}A+C_d\sigma}.
\]
This proves the suffix reduction with no change of start.

\paragraph{Prefix: Fourier conversion and low shells.}
Let \(\theta=\theta_i^{X,Y}(a,t)\), and let \(\zeta\) be the positive
distance pushforward of
\[
 M(a)^{-2}(D(x,y)2^a)^{-(d-1)/2}
 \prod_{\mathfrak t\in\mathcal M(a,t)}
 w_{i,\mathfrak t}(x,e(x,y))w_{i,\mathfrak t}(y,e(x,y))
 \,d\rho|_X(x)d\rho|_Y(y).
\]
On the good pairs defining \(\theta\), every mask is one.  The
remaining distance weight is bounded above and below by positive
dimensional constants on the separated geometry.  Thus positivity and
\eqref{rec:collision-fourier} imply
\begin{equation}\label{rec:prefix-lowpass}
 \mathsf D_i(a,t)
 \lesssim 2^{-a}\int_{|r|\le2^t}|\widehat\zeta(r)|^2\,dr.
\end{equation}
Since \(\zeta(\mathbb R)\le R^{C_d\sigma}\), frequencies below
\(2^a\) contribute \(R^{C_d\sigma}\), a stopping term.  Reality
allows us to treat only positive frequencies above \(2^a\), split
into shells \([2^{v-1},2^v]\), \(a<v\le t\).

Suppose \(t-v\ge\sigma N\).  A smooth majorant of the shell has
inverse Fourier transform bounded in modulus by
\(C2^v(1+2^v|s-s'|)^{-2}\).  Its contribution to
\eqref{rec:prefix-lowpass} is therefore at most
\[
 C2^{v-a}M(a)^{-4}
 \iint(1+2^v|D(x,y)-D(x',y')|)^{-2}
 \,d\theta'(x,y)d\theta'(x',y'),
\]
where \(\theta'\) is the unweighted pair restriction to the supports
of the old masks.  Every such pair satisfies
\(\mathcal M(a,v)\) at level \(i+1\), by list containment and mask
support.  Lemma~\ref{rec:binning} bounds the displayed expression by
\(C\mathsf D_{i+1}(a,v)\).  The normalization uses the same exact
\(M(a)\) on both sides; in particular, this step incurs no
\(R^{C_d\sigma}\) mass-comparison loss.  Summing shell indices costs
only \(R^{o(1)}\).

\paragraph{Prefix: high shells.}
Now suppose \(t-v<\sigma N\).  The largest nominal length in the
list is at most \((t-a)/2\).  Nonterminality gives
\[
 v-a>100q_0-\sigma N>20\sigma N,\qquad
 L\le(v-a)/2+\sigma N.
\]
Thus Proposition~\ref{osc:shell-conversion} applies to \(\zeta\).
For the prefix choice \(p\), the descriptor \([a,p]\) is present,
and
\begin{align*}
 p-a&\le t-p\le v-p+\sigma N,\\
 L&\le(t-a)/2\le t-p\le v-p+\sigma N,\\
 v-p&\ge(t-a)/2-\sigma N>10\sigma N.
\end{align*}
These verify the hypotheses of Corollary~\ref{graph:separated-form},
with \(g=a\).  Its one-profile specialization first gives general
independent energies on depth-\(p\) cells, formed with the restricted
converted symbols and still carrying every old mask.  In those positive
energies, remove the factors whose anchors are at most \(p\).  The
remaining list at each endpoint is exactly
\(\mathcal M_{>p}([a,t])\), and
\(a\le p\le v\le t\le v+\sigma N\).  The resulting values are therefore
recursive states \(\mathsf H_i^{[a,t]}(p,v)\), and the high contribution
is bounded by
\[
 R^{-2(A_p-A_a)+k(a,p)+C_d\sigma+o(1)}
       \sup\mathsf H_i^{[a,t]}(p,v)+O(R^{-T}).
\]
Conversion produced normalized angular derivatives of masks at index
\(i\), with a fixed enlarged budget and common angular bounds.  Their
supports still impose the existing level-\(i+1\) tests, which are accepted
by the graph estimate.  Thus the new angular state keeps index \(i\).

\paragraph{Direct angular reduction.}
Consider \(\mathsf H_i^{\mathcal B}(a,v)\) with base
\(\mathcal B=[a_0,t_0]\).  Its remaining
anchors are strictly deeper than \(a\), and
\(t_0\le v+\sigma N\).  If \(h\le a\), nonterminality forces
\(L=L_0\), since otherwise
\(v-a\le2\lceil10\sigma N\rceil<100q_0\).
Choose a maximal-length descriptor, with anchor \(p>a\) and an
origin \([b_0,u_0]\) in the base.  Its origin satisfies
\[
 b_0\le p-L,\qquad p+L\le u_0\le v+\sigma N,
 \qquad v-a\le2L.
\]
It follows that
\[
 L\le v-p+\sigma N,\qquad
 p-a\le v-p+2\sigma N,\qquad
 v-p\ge(v-a)/2-\sigma N>10\sigma N.
\]
Set \(g=\max(a,p-L)\).  Then
\[
 a\le g<p,\qquad
 p-L\le v+\sigma N-2L\le a+\sigma N,
 \qquad g\le a+\sigma N.
\]
The same origin supplies the parallel descriptor \([g,p]\), whose
anchor has not yet been reached.  Proposition~\ref{graph:estimate}
therefore applies to the general energy underlying
\(\mathsf H_i^{\mathcal B}(a,v)\), without any separation of its two
cells, and gives the direct reduction and the cost \(k(a,p)\).
The general graph child still carries the old symbols.  Remove its
reached factors only after forming its positive independent energy.  The
remaining list is exactly \(\mathcal M_{>p}(\mathcal B)\), so this gives
\(\mathsf H_i^{\mathcal B}(p,v)\) with the same fixed symbol bounds.

\paragraph{Near angular reduction.}
If \(h>a\), split each transform into its depth-\(h\) cell pieces,
temporarily keeping \emph{all} remaining factors.  Expand the two
squares in Equation~\eqref{osc:H}.  A near term is one for
which both center displacements, within the first and the second
parent cell respectively, are at most \(10(d+1)2^{-h}\).
In the graph with row \((P,Q)\) and column \((P',Q')\), these
conditions give bounded row and column degrees: a fixed grid cell has
only a dimensional number of neighbors at this distance.  Apply
\(2|uv|\le|u|^2+|v|^2\) to each row--column product and integrate
the resulting independent-angle squares.  There are at most
\(R^{2I(a,h)+C_d\sigma}\) rows or columns.  The four transform
splitting factors, the label count, and the ratio of energy
prefactors give
\[
 2^{d(h-a)}R^{-4I(a,h)}R^{2I(a,h)+C_d\sigma}
 =R^{-2(A_h-A_a)+C_d\sigma},
\]
by \eqref{rec:normalization}.  This is the zero-cost near estimate.
Only after obtaining each positive independent energy do we remove factors
anchored at depths at most \(h\); on its spatial cells they are angular
factors of bounded modulus.  The remaining list is exactly
\(\mathcal M_{>h}(\mathcal B)\), giving the child
\(\mathsf H_i^{\mathcal B}(h,v)\) with the same fixed symbol bounds.
This removal does not apply to the far terms that follow.

\paragraph{Far angular reduction: analytic estimate.}
In a remaining term at least one point displacement has length
\(\gtrsim_d2^{-h}\).  Fix the four spatial points and denote the two
displacements by \(\Delta_1,\Delta_2\), with lengths \(d_1,d_2\).
Put \(d_+=\max(d_1,d_2)\), \(d_-=\min(d_1,d_2)\).
For the larger displacement, on the radial support,
\[
 \Lambda=r d_+\gtrsim 2^{v-h}=2^{2L_0},\qquad
 \frac{D_*^2}{\Lambda}
 \lesssim
 \begin{cases}
 R^{-18\sigma},&2^LR^{-50\sigma}\le1,\\
 R^{-98\sigma},&2^LR^{-50\sigma}>1.
 \end{cases}
\]
Also \(\Lambda\ge R^{c\sigma}\) and \(D_*\le R^2\) for
sufficiently large \(N\).  Lemma~\ref{osc:stationary} therefore
expands the corresponding full-sphere integral to any prescribed
negative power accuracy.  After extracting the leading stationary
gain, the coefficients and any fixed number of derivatives in
\(r/2^v\) are bounded: the factor \((r d_+)^{-j}\) pays for the
order-\(2j\) symbol derivatives.

If \(d_+>2d_-\), leave the other angular integral exact.  The
derivative in \(r\) of each resulting phase has modulus at least
\(d_+-d_->d_+/2\), uniformly in the unexpanded angle.  Repeated
integration by parts in \(r/2^v\) gives \(O(R^{-T})\), including
when \(d_-=0\).  It is legitimate to leave that integral exact
because its symbols are independent of \(r\).

Otherwise the lengths are comparable, and the second angular
integral admits the same expansion.  Choose an integer \(n\) with
\(d_+\in(2^{-n-1},2^{-n}]\), and set
\(b=n+\log_2K_d\).  Then both lengths lie between
\(K_d2^{-b}/4\) and \(K_d2^{-b}\).  Moving to their depth-\(b\)
cell centers changes these bounds by at most
\(2\sqrt d\,2^{-b}\); the choice of \(K_d\) puts both pairs in
the exact separated geometry of Definition~\ref{osc:positive-count}.
Since each displacement joins points in one depth-\(a\) cell,
\(d_+\le\sqrt d\,2^{-a}\), and hence \(b>a\).  The far lower
bound gives \(b\le h+O_d(1)<v\).

The two leading stationary gains multiply to
\(O_d(2^{-(d-1)(v-b)})\).  Radial integration, using the coefficient
derivative bounds just established, bounds each remaining kernel by
\begin{equation}\label{rec:far-kernel}
 C2^{-(d-1)(v-b)}
       (1+2^v|d_1-d_2|)^{-2}.
\end{equation}
For opposite stationary signs the phase involves \(d_1-d_2\);
for equal signs it involves \(d_1+d_2\), whose modulus is no
smaller.  This proves \eqref{rec:far-kernel} for every sign choice.

\paragraph{Far angular reduction: support and normalization.}
The interval \([b,v]\) is contained in the base, so
\(\mathcal M(b,v)\) is a sublist of its list.  Every anchor of this
child list is strictly deeper than \(b>a\), and no such factor has
been removed.  In particular none was removed at depth \(h\): the
far expansion used the symbols before the removal made for the near
child.  Each coefficient differentiates the original masks only
angularly, so it retains all of their supports.  For example, the
amplitude for the internal pair \((x,x')\) is
\(b_x(w)\overline{b_{x'}(w)}\).  Its stationary coefficient is evaluated
at \(w=\pm(x-x')/|x-x'|\); a nonzero coefficient therefore imposes the
level-\(i+1\) tests of \(\mathcal M(b,v)\) at both spatial endpoints
\(x\) and \(x'\).  The same assertion holds for the other internal pair
\((y,y')\).  Antipodal symmetry of the tests handles the negative
stationary direction.  The expansion remainders are already
arbitrary-power errors and require no support assertion.

The far terms were obtained by splitting the transforms at depth
\(h\).  Once their four cell transforms are opened, the four splitting
factors cancel their model denominators exactly:
\[
 \left(\frac{M(h)}{M(a)}\right)^4M(h)^{-4}=M(a)^{-4}.
\]
For each comparable length band defining \(b\), the stationary estimate
is now a nonnegative kernel bound with the child-support indicators.
On this upper bound we may discard the indicator selecting the far
\(h\)-cell labels, while keeping the comparable length band.  Summing the
disjoint \(h\)-cell partition reassembles the raw four-point measure on
\(X^2\times Y^2\), which we then partition into depth-\(b\) cells.
This explains why no count of \(h\)-cell labels remains, whether \(b\)
is smaller or larger than \(h\).

For a fixed quadruple of separated depth-\(b\) cells \((P,P',Q,Q')\),
put
\[
 \eta_1=D_\#\theta_{i+1}^{P,P'}(b,v),\qquad
 \eta_2=D_\#\theta_{i+1}^{Q,Q'}(b,v).
\]
These are the unnormalized positive distance measures for the two
internal pairs, each restricted by the child tests at both endpoints.
By Lemma~\ref{rec:binning} and the exact normalization of
\(\mathsf D_{i+1}\),
\[
 \iint(1+2^v|s-s'|)^{-2}\,d\eta_1(s)d\eta_2(s')
 \lesssim M(b)^4 2^{-(v-b)}
              \max\mathsf D_{i+1}(b,v).
\]
Indeed each bin-square sum is bounded by its unshifted positive
collision count, and the geometric mean of the two counts is at
most their maximum.  Once the first cell of either internal pair
is fixed, its partner has only a dimensional number of possible
labels, since their center distance is \(O_d(K_d2^{-b})\).
There are thus at most \(R^{2I(a,b)+C_d\sigma}\) quadruple labels.
The two stationary gains and the scalar collision normalization combine
to \(2^{-d(v-b)}\), converting the parent frequency prefactor into
\(2^{d(b-a)}\).  The four depth-\(b\) masses and the two free first-cell
labels leave \(R^{-2I(a,b)+C_d\sigma}\).  Explicitly,
\begin{align*}
 &2^{d(v-a)}M(a)^{-4}
   2^{-(d-1)(v-b)}
   M(b)^4 2^{-(v-b)}R^{2I(a,b)+C_d\sigma}\\
 &\hspace{20mm}
 =2^{d(b-a)}R^{-2I(a,b)+C_d\sigma}
 =R^{-2(A_b-A_a)+C_d\sigma}.
\end{align*}
Summing the possible scales \(b\), coefficient terms, and shell
indices introduces only \(R^{o(1)}\).  This proves the zero-cost
far reduction, with one increase of the smoothing index.

\paragraph{Stopping.}
For a terminal interval, its width is at most \(100q_0\).
The Definitions of the two state values, bounded symbol
moduli, and the occupied-cell bounds give respectively
\[
 \mathsf H_i^{\mathcal B}(a,v)\lesssim2^{d(v-a)}R^{C_d\sigma},\qquad
 \mathsf D_i(a,t)\lesssim2^{t-a}R^{C_d\sigma}.
\]
Since \(q_0/N=\kappa+O(1/N)\), both satisfy
\eqref{rec:terminal}.  The low-frequency stopping term was already
estimated after \eqref{rec:prefix-lowpass}.  This completes the proof.
\end{proof}
